\documentclass[main.tex]{subfiles}

\title{Dirichlet Integral}
\author{Senan Sekhon}
\date{\DTMToday}
% Created September 15, 2021

\begin{document}

\maketitle

\VersionTimestamp

\section{Introduction}

\begin{lemma}[Dirichlet Kernel]\label{dirichlet_kernel}
    For all $n\in\N_0$ and all $x\in\R$, we have:
    \begin{equation*}
        1+2\sumkn\cos(2kx) = \frac{\sin((2n+1)x)}{\sin(x)}
    \end{equation*}
\end{lemma}
\begin{proof}
    Clearly the identity holds for $n=0$ as both sides are $1$. Suppose the identity holds for some $n\in\N_0$. We want to show it holds for $n+1$. We have:
    \begin{align*}
        1+2\sum_{k=1}^{n+1}\cos(2kx)
        &= 1+2\sumkn\cos(2kx)+2\cos((2n+2)x)=\frac{\sin((2n+1)x)}{\sin(x)}+2\cos((2n+2)x) \\
        &= \frac{\sin((2n+1)x)+2\sin(x)\cos((2n+2)x)}{\sin(x)}
    \end{align*}
    Note that $\sin((2n+3)x)-\sin((2n+1)x)=2\sin(x)\cos((2n+2)x)$. This yields:
    \begin{equation*}
        1+2\sum_{k=1}^{n+1}\cos(2kx)
        = \frac{\sin((2n+1)x)+\sin((2n+3)x)-\sin((2n+1)x)}{\sin(x)}
        = \frac{\sin((2n+3)x)}{\sin(x)}
    \end{equation*}
    Thus the identity holds for $n+1$, and so it holds for all $n\in\N_0$.
\end{proof}

\begin{theorem}[Dirichlet Integral]
    \begin{equation*}
        \int_0^\infty \frac{\sin(x)}{x}\dx = \frac{\pi}{2}
    \end{equation*}
\end{theorem}
\begin{proof}
    Define $K:\qty[0,\frac{\pi}{2}]\to\R$ as follows:
    \begin{equation*}
        K(x) = \begin{cases}
            \dfrac{\sin((2n+1)x)}{\sin(x)} & x\ne 0 \\
            2n+1 & x=0
        \end{cases}
    \end{equation*}
    Clearly $K$ is continuous on $\big(0,\frac{\pi}{2}\big]$, and:
    \begin{align*}
        \lim_{x\to0^+} K(x)
        &= \lim_{x\to0^+} \frac{\sin((2n+1)x)}{\sin(x)} \\
        &= \lim_{x\to0^+} (2n+1)\frac{\sin((2n+1)x)}{(2n+1)x}\frac{x}{\sin(x)} \\
        &= (2n+1)\qty(\lim_{x\to0^+} \frac{\sin((2n+1)x)}{(2n+1)x})\qty(\lim_{x\to0^+} \frac{x}{\sin(x)}) \\
        &= (2n+1) \cdot 1 \cdot 1 \\
        &= 2n+1
    \end{align*}
    Thus $K$ is continuous (and thus Riemann integrable) on $\qty[0,\frac{\pi}{2}]$. Using the \hyperref[dirichlet_kernel]{Dirichlet kernel}, we have:
    \begin{equation*}
        \int_0^{\pi/2} \frac{\sin((2n+1)x)}{\sin(x)}\dx
        = \int_0^{\pi/2} \qty(1+2\sumkn\cos(2kx))\dx
        = \int_0^{\pi/2} 1\dx + 2\sumkn\int_0^{\pi/2} \cos(2kx)\dx
        = \frac{\pi}{2}
    \end{equation*}
    We also have:
    \begin{equation*}
        \int_0^{\pi/2} \frac{\sin((2n+1)x)}{x}\dx
        = \int_0^{(2n+1)\pi/2} \frac{\sin(x)}{x}\dx
    \end{equation*}
    As $n\to\infty$, this tends to $\int_0^\infty \frac{\sin(x)}{x}\dx$. Define $g:(0,\frac{\pi}{2}]\to\R$ by $g(x)=\frac{1}{x}-\frac{1}{\sin(x)}$. Clearly $g$ is continuous on $\left(0,\frac{\pi}{2}\right]$. Also, using L'Hôpital's rule twice, we have:
    \begin{equation*}
        \lim_{x\to0^+} g(x)
        = \lim_{x\to0^+} \frac{\sin(x)-x}{x\sin(x)}
        = \lim_{x\to0^+} \frac{\cos(x)-1}{\sin(x)+x\cos(x)}
        = \lim_{x\to0^+} \frac{-\sin(x)}{2\cos(x)-x\sin(x)}
        = \frac{0}{2-0}
        = 0
    \end{equation*}
    Thus $g$ can be made continuous on $\qty[0,\frac{\pi}{2}]$ by setting $g(0)=0$. By the Riemann-Lebesgue lemma, we have $\int_0^{\pi/2} \qty(\frac{1}{x}-\frac{1}{\sin(x)})\sin((2n+1)x)\dx\to 0$.
\end{proof}

\end{document}